\begin{essayonly}
Trong phần này, các ứng dụng từ phần trước đã được sử dụng để xây dựng các thuật toán xấp xỉ điểm chịu ảnh hưởng bởi sai số. Đáng chú ý, S. Salzo và S. Villa đã chỉ ra một lỗ hổng nhỏ trong chứng minh của Guler, vốn làm ảnh hưởng đến kết luận về tốc độ hội tụ của thuật toán xấp xỉ điểm đối với hàm bậc hai.
\\
Mặc dù S. Salzo và S. Villa đã khắc phục được vấn đề này và khôi phục lại tốc độ hội tụ của thuật toán khi xét đến nhiều dạng sai số, kết quả họ thu được chỉ dừng lại ở mức biến đổi tuyến tính—bất kể sai số có tiến về không nhanh hay chậm. Điều này gợi ý rằng việc sử dụng thuật toán tăng tốc với sai số loại 1 thực tế không mang lại ưu thế rõ rệt, vì hiệu năng của nó chỉ tương đương với phiên bản không tăng tốc.
\\
Để giải quyết hạn chế trên, họ đã tiếp tục nghiên cứu thuật toán dưới cơ chế sai số loại 2 (với các yêu cầu khắt khe hơn loại 1) và thành công đạt được tốc độ hội tụ bậc 2, miễn là sai số giảm đủ nhanh. Hơn nữa, ngay cả khi chỉ sử dụng sai số loại 2, thuật toán vẫn duy trì được tốc độ hội tụ bậc 2 kể cả khi sai số giảm chậm hơn tiêu chuẩn, với tổng sai số tiến tới vô cùng.
\end{essayonly}